%% ===========================================================================
%%  試験・小テスト用テンプレート（堀田グループ）
%%
%%  このテンプレートは MIT License です。自由に使ってください。
%%  https://github.com/hottahd/latex-template
%%
%%  エンジンは LuaLaTeX です。このディレクトリで
%%      latexmk
%%  とすれば main.pdf ができます。
%%
%%  問題用と解答例つきの 2 種類を、下の 1 行を書き換えるだけで作り分けられます。
%%  問題文は 1 つしかないので、両者が食い違う心配がありません。
%% ===========================================================================

\documentclass[a4j,12pt]{ltjarticle}
%\documentclass[a4j,twocolumn]{ltjarticle}   % 2 段組で詰めたいとき

%% ===========================================================================
%%  1. どちらを出力するか
%% ---------------------------------------------------------------------------
%%   \showanswerfalse  問題用。\begin{kaitou} ... \end{kaitou} は出力されない
%%   \showanswertrue   解答例つき。表題に「解答例」がつく
%%  作り分けたら PDF の名前を変えて保存すること。
%%      latexmk && mv main.pdf 期末テスト_問題.pdf
%% ===========================================================================
\newif\ifshowanswer
\showanswerfalse
%\showanswertrue


%% ===========================================================================
%%  2. パッケージと設定
%% ===========================================================================

%% 共通プリアンブル（フォント・数式・図・単位・マクロ）。
%% 中身は achilles.sty を見てください。全テンプレートで共通です。
\usepackage[nogeometry]{achilles}
\usepackage[margin=20truemm]{geometry}

\usepackage{comment}     % 解答を丸ごと出し入れするために使う
\usepackage{tcolorbox}
\tcbuselibrary{skins,breakable}

%% 節番号を「1」「2」ではなく「問 1」「問 2」にする
\renewcommand{\thesection}{問\arabic{section}}
\renewcommand{\thesubsection}{(\arabic{subsection})}
%% 式番号は節ごとに振り直す
\numberwithin{equation}{section}

%% --- 解答環境 -------------------------------------------------------------
%% 使い方（\end{kaitou} は必ず独立した行に書くこと。comment パッケージの制約）
%%
%%   \begin{kaitou}
%%     $T = 2\pi/\Omega = \qty{25}{day}$ である。
%%   \end{kaitou}
%%
\ifshowanswer
  \specialcomment{kaitou}{%
    \begin{tcolorbox}[breakable,colback=blue!3,colframe=blue!45,
                      title=解答,fonttitle=\bfseries\small,left=4pt,right=4pt]%
  }{%
    \end{tcolorbox}%
  }
\else
  \excludecomment{kaitou}
\fi

%% --- 解答欄 ---------------------------------------------------------------
%% \ansspace{4cm} で、問題用のときだけ 4cm の余白を空ける（答案に書かせる欄）。
%% 解答例つきのときは詰めて出す。
\newcommand{\ansspace}[1]{\ifshowanswer\vspace{2mm}\else\vspace{#1}\fi}

%% --- 配点 -----------------------------------------------------------------
%% 見出しの右端に（10 点）と出す。\section{太陽の自転\haiten{20}} のように、
%% 見出しの中に書く。
\newcommand{\haiten}[1]{\texorpdfstring{\hfill\normalsize（#1 点）}{}}  % しおりには出さない

%% --- 表題 -----------------------------------------------------------------
%% ここを書き換える
\newcommand{\kamoku}{宇宙物理学III}
\newcommand{\shiken}{期末テスト}
\newcommand{\jisshibi}{2026年7月30日}

\title{\kamoku\ \shiken\ifshowanswer\ 解答例\fi}
\author{}
\date{\jisshibi}


%% ===========================================================================
%%  3. 本文
%% ===========================================================================
\begin{document}
\maketitle
\thispagestyle{empty}

%% 問題用のときだけ、注意書きと氏名欄を出す
\ifshowanswer\else
\noindent
\begin{tcolorbox}[colback=black!3,colframe=black!40,left=6pt,right=6pt]
  \small
  \begin{itemize}[leftmargin=1.2em,itemsep=0pt,topsep=0pt]
    \item 解答は答案用紙に書くこと。途中の考え方も採点の対象とする。
    \item 電卓・参考書の持ち込みは不可。
    \item 有効数字の指定がある問いは、指定の桁で答えること。
  \end{itemize}
  \vspace{2mm}
  \normalsize
  学籍番号 \underline{\hspace{40mm}}\qquad 氏名 \underline{\hspace{50mm}}
\end{tcolorbox}
\vspace{5mm}
\fi


\section{太陽の自転\haiten{20}}

太陽の赤道は $\Omega/2\pi = \qty{460}{nHz}$ 程度で自転している。
ここで $\Omega = u_\phi/L$ は角速度、$u_\phi$ は自転速度、$L$ は回転軸からの距離である。
太陽半径は \qty{700}{Mm} とする。

\subsection{}
赤道の自転周期を有効数字 2 桁で求めなさい（単位は日）。

\ansspace{3cm}

\begin{kaitou}
  $T = 1/(\Omega/2\pi) = 1/(\qty{460e-9}{\per\second})
   = \qty{2.17e6}{\second} \simeq \qty{25}{\day}$ である。
\end{kaitou}

\subsection{}
赤道の自転速度を有効数字 1 桁で求めなさい（単位は $\mathrm{km~s^{-1}}$）。

\ansspace{3cm}

\begin{kaitou}
  $u_\phi = 2\pi R_\odot / T
   = 2\pi \times \qty{700}{Mm} / \qty{2.17e6}{\second}
   \simeq \qty{2}{\km\per\second}$ である。
\end{kaitou}


\section{見かけの自転周期\haiten{15}}

地球は公転しているので、地球から観測した見かけの太陽の自転周期は
問 1 の答えよりもやや長くなる。どの程度長くなるか有効数字 1 桁で答えよ（単位は日）。
地球の公転角速度を $\Omega_\oplus$ として、$\Omega_\oplus/\Omega$ の
一次までの寄与を考えればよい。

\ansspace{5cm}

\begin{kaitou}
  見かけの角速度は $\Omega - \Omega_\oplus$ なので、
  \begin{align*}
    T_\mathrm{syn} = \frac{2\pi}{\Omega-\Omega_\oplus}
      = \frac{T}{1-\Omega_\oplus/\Omega}
      \simeq T\left(1+\frac{\Omega_\oplus}{\Omega}\right).
  \end{align*}
  $\Omega_\oplus/\Omega = T/T_\oplus = 25/365 \simeq 0.068$ より、
  延びは $T \times 0.068 \simeq \qty{2}{\day}$ である。
\end{kaitou}

\end{document}
